% Generated draft: requires mathematical and source review.
\documentclass[11pt]{article}
\usepackage[margin=25mm]{geometry}
\usepackage{fontspec}
\setmainfont{Latin Modern Roman}
\usepackage{amsmath,amssymb,mathtools,unicode-math}
\setmathfont{Latin Modern Math}
\usepackage{enumitem,needspace,xurl,hyperref}
\setlength{\parindent}{0pt}
\setlength{\parskip}{5pt}
\setlength{\emergencystretch}{3em}
\newcommand{\N}{\mathbb{N}}
\newcommand{\Z}{\mathbb{Z}}
\newcommand{\Q}{\mathbb{Q}}
\newcommand{\R}{\mathbb{R}}
\newcommand{\C}{\mathbb{C}}
\newcommand{\F}{\mathbb{F}}
\newcommand{\A}{\mathbb{A}}
\newcommand{\PP}{\mathbb{P}}
\newcommand{\E}{\mathbb{E}}
\newcommand{\eps}{\varepsilon}
\newcommand{\dd}{\,\mathrm{d}}
\begin{document}
\section*{UnsolvedMath ID: 701026}
\textbf{Source problem code:} catalogue entry
\subsection{Definitions and mathematical statement}
For a group \(G\), define its lower central series by
\[
\gamma_1(G)=G,\qquad \gamma_{i+1}(G)=[\gamma_i(G),G]\quad(i\geq 1),
\]
where \([H,G]\) is the subgroup generated by the commutators \(h^{-1}g^{-1}hg\) with \(h\in H\) and \(g\in G\). Following the supplied definition, call \(G\) \emph{para-free} if there exists a free group \(F\) such that
\[
\gamma_i(G)/\gamma_{i+1}(G)\cong
\gamma_i(F)/\gamma_{i+1}(F)\quad\text{for every }i\geq 1,
\qquad
\bigcap_{i=1}^{\infty}\gamma_i(G)=\{1\}.
\]
No finite generation or restriction on the rank of \(F\) is imposed. The displayed isomorphisms are isomorphisms of groups; the supplied definition does not require compatibility between them or isomorphisms of the corresponding nilpotent quotients.

The question of A. N. Zubkov is whether every such group \(G\) is representable by matrices over a commutative associative ring with identity. Interpreting matrix representability in its usual faithful, finite-dimensional sense, the question is whether, for every para-free group \(G\), there exist a positive integer \(n\), a commutative associative ring \(R\) with identity, and an injective group homomorphism
\[
\rho:G\hookrightarrow\operatorname{GL}_n(R),
\]
where \(\operatorname{GL}_n(R)\) denotes the group of invertible \(n\times n\) matrices over \(R\). Both \(n\) and \(R\) may depend on \(G\). The source does not explicitly specify faithfulness or finite matrix size; these are the interpretive qualifications in this formulation.
\subsection{Short English statement}
Does every para-free group, as defined by its lower central factors and trivial intersection of its lower central series, admit a faithful representation by finite square matrices over some commutative associative ring with identity? The supplied record labels this question open; that status has not been independently verified.
\subsection{Sources}
\begin{itemize}
\item Local UnsolvedMath registry: permanent ID \(701026\), exact source code \texttt{catalogue entry}, title \emph{Kourovka Notebook Problem 14.41}. The supplied record attributes the problem to A. N. Zubkov, with proposed year 1999, and labels it open.
\item Supplied source attribution: \emph{The Kourovka Notebook (Unsolved Problems in Group Theory)}, original problem 14.41, 14th issue (1999); also identified in the record with the 21st issue edition (2026), arXiv:1401.0300v48. The supplied collection description names E. I. Khukhro and V. D. Mazurov as editors. This bibliographic attribution is supplied evidence and has not been independently checked.
\end{itemize}
\end{document}
